\documentclass[11pt]{article}
\usepackage[T1]{fontenc}
\usepackage{lmodern}
\usepackage[margin=1.05in]{geometry}
\usepackage{amsmath,amssymb,amsthm,mathtools}
\usepackage{microtype,booktabs,array}
\usepackage{xcolor}
\usepackage[colorlinks=true,linkcolor=blue!45!black,citecolor=blue!45!black,urlcolor=blue!45!black]{hyperref}
\usepackage[nameinlink,capitalise,noabbrev]{cleveref}
\usepackage{newunicodechar}
\newunicodechar{χ}{\ensuremath{\chi}}
\newunicodechar{∃}{\ensuremath{\exists}}
\newunicodechar{∀}{\ensuremath{\forall}}
\newunicodechar{∧}{\ensuremath{\wedge}}
\newunicodechar{∨}{\ensuremath{\vee}}
\newunicodechar{¬}{\ensuremath{\neg}}
\newunicodechar{→}{\ensuremath{\to}}
\newunicodechar{≤}{\ensuremath{\le}}
\newunicodechar{≥}{\ensuremath{\ge}}
\newunicodechar{∈}{\ensuremath{\in}}
\newunicodechar{₅}{\ensuremath{_5}}
\newtheorem{theorem}{Theorem}[section]
\newtheorem{proposition}[theorem]{Proposition}
\newtheorem{lemma}[theorem]{Lemma}
\newtheorem{corollary}[theorem]{Corollary}
\theoremstyle{definition}
\newtheorem{definition}[theorem]{Definition}
\theoremstyle{remark}
\newtheorem{remark}[theorem]{Remark}
\newcommand{\F}{\mathbb F}
\newcommand{\GF}{\mathrm{GF}}
\newcommand{\AGL}{\mathrm{AGL}}
\newcommand{\RFSF}{98}
\newcommand{\RFSFq}{97}
\newcommand{\RFSFgroup}{$\AGL(1,97)$}
\newcommand{\RFSFtext}{$q=97$ and $C=\F_{97}^{*}$: the colouring of the $4$-subsets of $\F_{97}$ is invariant
under the full affine group $\AGL(1,97)$, so it is determined by the colours of the sets $\{0,1,t,u\}$; the
$1752$ red pairs $\{t,u\}$ (out of $4465$) are listed in \texttt{data/describe\_R564\_97.txt}. There are $385$ orbits
on $4$-sets and $1\,359\,552$ of the $3\,464\,840$ quadruples are red. There is no red $K_5^{(4)}$ and no blue
$K_6^{(4)}$. (Colourings with the same symmetry type exist for every prime $67\le q\le 97$ that we tried.)}
\newcommand{\RLeanTimes}{on a Xeon Gold 6254 node, Lean 4.35.0, one core:
the whole file \texttt{lean/Main.lean} needs $4595$\,s CPU ($1$\,h\,$21$\,min wall) and $13.7$\,GB, most of it
for \texttt{R564} (individually, on another machine: \texttt{R453} $5$\,s, \texttt{R453f} $53$\,s,
\texttt{R553} $331$\,s, \texttt{R463} $128$\,s, \texttt{R554} $26$\,s); the full output, including the
\texttt{\#print axioms} lines, is in \texttt{lean/lean\_output.txt}}
\newcommand{\RSweepFive}{Smaller $C$ ($m\in\{25,34\}$) at $q=101,103$ and the one-vertex extension of the $97$-point colouring (a SAT problem with $4656$ variables) were not decided within the time limits.}
\newcommand{\RSweepFourFive}{Not decided within the time limits: $q=43,49,61,67$ with $m=3$, the cyclic group on $38$ points, $\mathbb Z_{19}\rtimes C_3$ acting on two copies of $\F_{19}$ ($38$ points), and the one-vertex extension of the $37$-point colouring ($666$ variables). The extension of the $37$-point family by a fixed point ($G(37,C_3)$ acting on $38$ points) is UNSAT, as are its analogues with $2$ and $3$ fixed points and with $|C|\in\{6,9\}$.}
\newcommand{\RSweepFourSix}{Also UNSAT for $q=71$ with $m\in\{10,14\}$ and $q=73$ with $m\in\{8,12\}$.}
\newcommand{\RSweepFFour}{For $R(5,6;4)$ every prime $67\le q\le 103$ tried was SAT.}
\title{New lower bounds for hypergraph Ramsey numbers\\ from affine-invariant colourings}
\author{Claude (AI model by Anthropic)\\ \small run and submitted by Aleksandr Bersenev}
\date{October 14, 2026\\[3pt]\footnotesize Produced autonomously by Claude (Opus 5.5) agents in Claude Code; run and submitted by Aleksandr Bersenev.\\ \footnotesize Not an official Anthropic publication. Not peer reviewed.}
\hypersetup{pdftitle={New lower bounds for hypergraph Ramsey numbers from affine-invariant colourings},pdfauthor={Claude (AI model by Anthropic); submitted by Aleksandr Bersenev}}
\begin{document}
\maketitle
\begin{abstract}
We prove
\[
  R(4,5;3)\ge 38,\quad R(5,5;3)\ge 98,\quad R(4,6;3)\ge 68,\quad R(5,5;4)\ge 38,\quad R(5,6;4)\ge \RFSF,
\]
improving the lower bounds $35$, $88$, $63$, $35$ and $67$ listed in revision~\#18 (April 2026) of
Radziszowski's dynamic survey \emph{Small Ramsey Numbers}. Each bound is witnessed by an explicit
red/blue colouring of the $r$-subsets of a prime field $\GF(q)$ that is invariant under a group of affine
maps $x\mapsto ax+b$, where $a$ ranges over a fixed subgroup of $\F_q^{*}$ (the cube roots of unity for
$R(4,5;3)$, the nonzero squares for $R(5,5;3)$ and $R(4,6;3)$, all of $\F_q^{*}$ for the $4$-uniform cases). Imposing
the affine symmetry shrinks the satisfiability problem to a few dozen to a few hundred Boolean
variables, which a SAT solver handles in seconds. All five colourings are certified in Lean~4: the
absence of monochromatic cliques through a fixed vertex is checked by the Lean kernel, and a short
formal argument reduces arbitrary cliques to that case.
\end{abstract}

\section{Introduction}

For integers $k,l\ge r\ge 2$ the hypergraph Ramsey number $R(k,l;r)$ is the least $n$ such that every
colouring of the $r$-element subsets of an $n$-element set with the colours red and blue contains $k$
points all of whose $r$-subsets are red, or $l$ points all of whose $r$-subsets are blue. Such sets of
points are called a red $K_k^{(r)}$ and a blue $K_l^{(r)}$. For $r\ge 3$ essentially no values are known:
the only nontrivial classical value is $R(4,4;3)=13$~\cite{MR1}, and Section~7.1 of the
survey~\cite{DS1} lists only lower bounds for the next cases,
\[
 35\le R(4,5;3),\quad 63\le R(4,6;3),\quad 88\le R(5,5;3),\quad 79\le R(4,4,4;3),\quad 35\le R(5,5;4),
\]
together with bounds such as $R(5,6;4)\ge 67$ obtained from general recurrences
(\cite[7.1(f),(j)]{DS1}, \cite{Liu}). The bounds for $R(4,5;3)$, $R(4,6;3)$, $R(5,5;3)$ and $R(4,4,4;3)$ are due to
Dybizbański~\cite{Dyb2,Dyb3}, who found \emph{cyclic} (shift-invariant) colourings with a SAT solver;
before that $R(5,5;3)\ge 82$ was known~\cite{BudHR}. The bound $R(5,5;4)\ge 35$ is due to Exoo~\cite{DS1}.

In 2026 an unrefereed public project~\cite{SV} used group-invariant colourings and SAT to obtain
$R(4,4,4;3)\ge 84$, $R(4,6;3)\ge 64$ and $R(5,5;4)\ge 36$. For $R(5,5;3)$ that project lists experiments
only at $n=88$, and an issue in its tracker (September 2026, self-reported and not reviewed) claims an
$88$-point colouring, i.e.\ $R(5,5;3)\ge 89$.

\begin{theorem}\label{thm:main}
$R(4,5;3)\ge 38$, $R(5,5;3)\ge 98$, $R(4,6;3)\ge 68$, $R(5,5;4)\ge 38$ and $R(5,6;4)\ge \RFSF$.
\end{theorem}

To our knowledge (literature and web search of October 14, 2026) all five bounds are new, and each improves
every bound we are aware of by at least two. The case $R(4,5;3)$ is the smallest open classical hypergraph
Ramsey number; its first lower bound, $24$, is due to Isbell, and the bound $33$ was improved to $35$ by
Dybizbański in 2018~\cite{Dyb2,DS1}. The proof is a construction: for a prime $q$ we colour
the $r$-subsets of $\GF(q)$ invariantly under a group of affine maps. \Cref{tab:results} summarises
the witnesses.

\begin{table}[h]
\centering
\begin{tabular}{lccccc}
\toprule
number & DS1.18 & other 2026 claims & this paper & $q$ & symmetry group\\
\midrule
$R(4,5;3)$ & $\ge 35$ & --- & $\ge 38$ & $37$ & $x\mapsto ax+b$, $a^3=1$\\
$R(5,5;3)$ & $\ge 88$ & $\ge 89$ (unreviewed) & $\ge 98$ & $97$ & $x\mapsto ax+b$, $a\in(\F_{97}^{*})^{2}$\\
$R(4,6;3)$ & $\ge 63$ & $\ge 64$ & $\ge 68$ & $67$ & $x\mapsto ax+b$, $a\in(\F_{67}^{*})^{2}$\\
$R(5,5;4)$ & $\ge 35$ & $\ge 36$ & $\ge 38$ & $37$ & $\AGL(1,37)$\\
$R(5,6;4)$ & $\ge 67$ & --- & $\ge \RFSF$ & $\RFSFq$ & \RFSFgroup\\
\bottomrule
\end{tabular}
\caption{The improved bounds. ``Other 2026 claims'' refers to~\cite{SV}.}
\label{tab:results}
\end{table}

\section{Affine-invariant colourings}

Let $q$ be a prime and let $C\le \F_q^{*}$ be a subgroup of order $m$. The group
$G=G(q,C)=\{x\mapsto ax+b:\ a\in C,\ b\in\F_q\}$ acts on $\F_q$, transitively (already the translations
do), and hence on the $r$-subsets of $\F_q$. A colouring $\chi$ of the $r$-subsets is \emph{$G$-invariant}
if $\chi(gS)=\chi(S)$ for all $g\in G$; such colourings correspond to colourings of the $G$-orbits of
$r$-subsets, of which there are roughly $\binom{q}{r}/(qm)$.

\begin{lemma}\label{lem:through0}
Let $\chi$ be invariant under all translations $x\mapsto x+b$ of $\F_q$. If no $k$-subset containing $0$ is
monochromatic, then no $k$-subset is monochromatic.
\end{lemma}
\begin{proof}
If $S$ is monochromatic and $s\in S$, then $S-s$ contains $0$ and is monochromatic of the same colour,
because $\chi(T-s)=\chi(T)$ for every $r$-subset $T\subseteq S$.
\end{proof}

Consequently the existence of a $G$-invariant colouring without red $K_k^{(r)}$ and without blue
$K_l^{(r)}$ is equivalent to the satisfiability of the CNF formula with one Boolean variable $x_O$ per orbit
$O$ (true = red) and the clauses
\[
 \bigvee_{T\in\binom{S}{r}} \neg x_{O(T)}\quad (S\ni 0,\ |S|=k),\qquad\qquad
 \bigvee_{T\in\binom{S}{r}} x_{O(T)}\quad (S\ni 0,\ |S|=l).
\]
Because the stabiliser of $0$ in $G$ (multiplication by $C$) permutes these clauses, many of them coincide;
after removing duplicates, the formulas for $m=(q-1)/2$ have only $30$--$40$ variables and at most a few times
$10^4$ clauses (for instance $34$ variables for $q=97$), and are decided by the SAT solver
\texttt{kissat}~\cite{kissat} within seconds. We swept primes (and a few prime powers) $q$ above the
known bounds and divisors $m$ of $q-1$ from large to small; \Cref{sec:search} lists what was tried.

\section{The constructions}\label{sec:constr}

Fix a prime $q$, the subgroup $C=(\F_q^{*})^2$ of nonzero squares and a non-square $g$. If $\chi$ is a
$G(q,C)$-invariant colouring of the triples of $\F_q$, then for any labelling $\{x,y,z\}$ of a triple,
\[
  \chi\{x,y,z\}=\begin{cases}\chi\{0,1,t\}, & t=\frac{z-x}{y-x}, \text{ if } y-x \text{ is a square},\\[2pt]
  \chi\{0,g,t\}, & t=g\,\frac{z-x}{y-x}, \text{ if } y-x \text{ is a non-square}.\end{cases}
\]
Hence $\chi$ is determined by the two sets
$R_1=\{t:\chi\{0,1,t\}=\text{red}\}$ and $R_g=\{t:\chi\{0,g,t\}=\text{red}\}$.

\paragraph{$R(4,5;3)\ge 38$.} $q=37$ and $C=\{1,10,26\}$, the cube roots of unity in $\F_{37}$
($|G|=111$, $78$ orbits on triples), so the colouring is invariant under $x\mapsto x+1$ and $x\mapsto 10x$. Since $C$ is small we describe the colouring by the colours of the
triples $\{0,y,z\}$, $0<y<z<37$, which determine everything by translation: $\{x,y,z\}$ ($x<y<z$) has the
colour of $\{0,y-x,z-x\}$. Exactly $258$ of these $630$ triples are red; they are listed in
\texttt{data/describe\_R453\_37.txt} and encoded in the Lean file. In total $3182$ of the $7770$ triples are
red; there is no red $K_4^{(3)}$ and no blue $K_5^{(3)}$ (this was also checked by brute force over all
$\binom{37}{4}$ four-sets and $\binom{37}{5}$ five-sets).

\paragraph{$R(5,5;3)\ge 98$.} $q=97$, $g=5$,
\begin{align*}
R_1=\{&4, 5, 8, 9, 12, 13, 15, 20, 21, 22, 23, 24, 25, 26, 29, 31, 32, 35, 37, 38, 40, 41, 44, 54, 57, 58, 60,\\
 &61, 63, 66, 67, 69, 72, 73, 74, 75, 76, 77, 78, 83, 85, 86, 89, 90, 93, 94\},\\
R_5=\{&1, 4, 8, 9, 12, 13, 14, 16, 17, 18, 20, 24, 26, 27, 28, 29, 30, 31, 37, 38, 39, 40, 41, 42, 45, 47, 55,\\
 &57, 60, 61, 62, 63, 64, 65, 71, 72, 73, 74, 75, 76, 78, 82, 84, 85, 86, 88, 89, 90, 93, 94\}.
\end{align*}
The colouring has $34$ orbits on triples; $74\,496$ of the $147\,440$ triples are red. No five points of
$\F_{97}$ span a monochromatic $K_5^{(3)}$. (An analogous colouring of $\F_{89}$, which already gives
$R(5,5;3)\ge 90$, is in the data files.)

\paragraph{$R(4,6;3)\ge 68$.} $q=67$, $g=2$,
\begin{align*}
R_1&=\{2, 4, 7, 9, 11, 13, 15, 16, 17, 20, 21, 29, 30, 34, 37, 43, 46, 48, 57, 58, 61, 66\},\\
R_2&=\{1, 4, 9, 11, 14, 20, 22, 27, 29, 35, 37, 39, 40, 43, 44, 47, 50, 51, 55, 61, 62, 65\}.
\end{align*}
There is no red $K_4^{(3)}$ and no blue $K_6^{(3)}$; $16\,214$ of the $47\,905$ triples are red. Since
$67\equiv 3\pmod 4$, $-1$ is not a square and the group has odd order $67\cdot 33$.

\paragraph{$R(5,5;4)\ge 38$.} $q=37$ and $C=\F_{37}^{*}$, i.e.\ the colouring of the $4$-subsets of
$\F_{37}$ is invariant under the full affine group $\AGL(1,37)$, which acts sharply $2$-transitively;
every $4$-subset is equivalent to one of the form $\{0,1,t,u\}$, and the colouring is given by the list of
the $309$ red pairs $\{t,u\}\subseteq\F_{37}\setminus\{0,1\}$ in \texttt{data/describe\_R554\_37.txt}
($55$ orbits; $34\,299$ of the $66\,045$ quadruples are red). No five points span a monochromatic
$K_5^{(4)}$.

\paragraph{$R(5,6;4)\ge\RFSF$.} \RFSFtext

\section{Verification}\label{sec:verify}

Each colouring was checked in three independent ways.

\begin{enumerate}
\item A C program (\texttt{verify.c}, \texttt{verify4.c}) reads the list of red $r$-sets and searches, by
  depth-first clique extension over all vertices, for monochromatic $K_k^{(r)}$ of either colour; it reports
  none. It reports violations for perturbed inputs.
\item For the $3$-uniform colourings a Python script checks translation invariance and all
  $k$-subsets containing $0$; for $R(4,5;3)$ and $R(5,5;4)$ all $k$-subsets of the $37$ points were checked
  by brute force (e.g.\ all $\binom{37}{5}=435\,897$ five-sets).
\item A Lean~4 proof (below).
\end{enumerate}

\paragraph{Formal statement.} The Lean file \texttt{lean/Main.lean} (Lean~4.35.0, core library only, no
Mathlib) represents a colouring of the $3$-subsets of $\{0,\dots,N-1\}$ by a Boolean function on increasing
triples (\texttt{true} = red) and defines, e.g.,
{\small\begin{verbatim}
def HasMonoK4 (N : Nat) (col : Nat → Nat → Nat → Bool) (χ : Bool) : Prop :=
  ∃ a b c d : Nat, a < b ∧ b < c ∧ c < d ∧ d < N ∧
    col a b c = χ ∧ col a b d = χ ∧ col a c d = χ ∧ col b c d = χ
def HasMonoK5 (N : Nat) (col : Nat → Nat → Nat → Bool) (χ : Bool) : Prop :=
  ∃ a b c d e : Nat, a < b ∧ b < c ∧ c < d ∧ d < e ∧ e < N ∧
    col a b c = χ ∧ col a b d = χ ∧ col a b e = χ ∧ col a c d = χ ∧ col a c e = χ ∧
    col a d e = χ ∧ col b c d = χ ∧ col b c e = χ ∧ col b d e = χ ∧ col c d e = χ
\end{verbatim}}
and proves
{\small\begin{verbatim}
theorem R453.main : ∃ col : Nat → Nat → Nat → Bool,
    ¬ HasMonoK4 37 col true ∧ ¬ HasMonoK5 37 col false
theorem R553.main : ∃ col : Nat → Nat → Nat → Bool,
    ¬ HasMonoK5 97 col true ∧ ¬ HasMonoK5 97 col false
\end{verbatim}}
which say exactly that $R(4,5;3)>37$ and $R(5,5;3)>97$. The theorems \texttt{R463.main} ($N=67$, no red
$K_4^{(3)}$, no blue $K_6^{(3)}$), \texttt{R554.main} and \texttt{R564.main} (colourings
\texttt{col : Nat → Nat → Nat → Nat → Bool} of increasing quadruples, $5$ resp.\ $15$ quadruples per clique)
are analogous, and \texttt{R453f.main} re-proves $R(4,5;3)>37$ by a kernel search over all increasing
tuples that does not use any symmetry. A second file, \texttt{lean/MainFinset.lean} (with Mathlib), restates the two
$3$-uniform diagonal-type results with \texttt{Finset}s, e.g.
{\small\begin{verbatim}
theorem R453.main_finset : ∃ c : Finset (Fin 37) → Bool,
    (∀ S : Finset (Fin 37), S.card = 4 → ∃ t ∈ S.powersetCard 3, c t = false) ∧
    (∀ T : Finset (Fin 37), T.card = 5 → ∃ t ∈ T.powersetCard 3, c t = true)
\end{verbatim}}

\paragraph{Proof structure.} The colouring is \emph{defined} through differences,
\texttt{col x y z := Fmask.testBit ((y-x)*q + (z-x))}, where the natural number \texttt{Fmask} encodes the
colours of the triples $\{0,d_1,d_2\}$; therefore \texttt{col (x-a) (y-a) (z-a) = col x y z} for
$a\le x<y<z$ (lemma \texttt{col\_shift}, proved by \texttt{omega}), which is the formal version of
Lemma~\ref{lem:through0} with $s=\min S$ (no wrap-around is needed). The kernel then evaluates, by
\texttt{decide +kernel}, a Boolean function performing a pruned search for monochromatic cliques through
$0$: a vertex is only added if all new $r$-sets have the required colour. Two short lemmas
(\texttt{noRed0}, \texttt{noBlue0}) turn the value \texttt{true} of this search into the statement that no
monochromatic clique through $0$ exists. In the $4$-uniform files the colouring is
\texttt{col x y z w := (row (y-x)).testBit ((z-x)*q + (w-x))}, where \texttt{row} is a binary decision
tree returning masks of $q^2$ bits (this keeps the numbers handled by the kernel small), and the search is
split into one declaration per second vertex. \texttt{\#print axioms} reports only \texttt{propext} and
\texttt{Quot.sound} for all five main theorems; no \texttt{native\_decide} is used. Kernel times: \RLeanTimes.

\section{Searches that did not succeed}\label{sec:search}
For completeness we record the parameter ranges in which no colouring was found ($m=|C|$) (``UNSAT'' means the
SAT solver proved that no colouring with the stated symmetry exists).
\begin{itemize}
\item $R(5,5;3)$, $C=$ squares: SAT for $q=89, 97$; UNSAT for $q=101,103,107,109,113,121,125$.
\RSweepFive
\item $R(4,6;3)$, $C=$ squares: SAT for $q=67$; UNSAT for $q=71,73,79,81,83,89,97$. \RSweepFourSix
\item $R(5,5;4)$: SAT for $q=37$ ($m=36,18$); UNSAT for $q=41$ ($m=40,20$). \RSweepFFour
\item $R(4,5;3)$: SAT for $q=37$, $m=3$; UNSAT for $q=37$ with $m\in\{2,4,6,9,12,18,36\}$, for
  $q=41$ with $m\in\{2,4,5,8,10,20,40\}$, for $q\in\{43,47,49,53\}$ with large $m$, and for dihedral
  colourings with $34\le n\le 37$. \RSweepFourFive
\item $R(4,4,4;3)$ (three colours): UNSAT for $C$ of index $2$ and $3$ for $89\le q\le 125$.
\end{itemize}
These negative results only concern the very symmetric families searched; they say nothing about
$R(k,l;r)$ itself.

\section{Remarks and limitations}
\begin{itemize}
\item The theorems are lower bounds only; no upper bounds are claimed. All verification is of explicit
 finite objects; novelty is claimed to the best of our knowledge as of October 14, 2026.
\item By a general product-type construction quoted in~\cite[7.1(a)]{DS1} (attributed there to~\cite{BrBH}),
 a $2$-colouring witnessing $R(5,5;3)>m$ yields $R_k(5;3)> m^{2^{k-2}}$; with $m=97$ this would give
 $R(5,5,5;3)\ge 9410$. We did not formalise this consequence.
\item The sweep is far from exhaustive: smaller subgroups $C$, non-prime orders, and groups with fixed
 points were explored only partially, and several instances were stopped by a time limit; larger bounds are
 plausible. For $R(5,6;4)$ every prime $67\le q\le 103$ that we tried admits an $\AGL(1,q)$-invariant
 colouring; the colourings for $q=101$ and $q=103$ (giving $R(5,6;4)\ge 104$) were verified by the C
 program only, not in Lean, because the kernel check grows like $q^5$.
\item We also tried, without success, to decide $R(K_4-t,K_5;3)\in\{14,15,16\}$ by SAT: local search never
 gets below six violated constraints on $14$ points (while $13$ points are easy), which suggests
 $R(K_4-t,K_5;3)=14$, but we have no proof.
\end{itemize}

\begin{thebibliography}{9}
\bibitem{DS1} S.~P. Radziszowski, \emph{Small Ramsey Numbers}, Electronic Journal of Combinatorics,
 Dynamic Survey DS1, revision \#18, April 24, 2026. \url{https://doi.org/10.37236/21}.
\bibitem{Dyb2} J.~Dybizbański, A lower bound on the hypergraph Ramsey number $R(4,5;3)$,
 \emph{Contributions to Discrete Mathematics} 13(2) (2018) 112--115.
\bibitem{Dyb3} J.~Dybizbański, addendum to~\cite{Dyb2} (colourings for $R(5,5;3)\ge 88$, $R(4,6;3)\ge 63$,
 $R(4,4,4;3)\ge 79$), \url{https://inf.ug.edu.pl/ramsey}, as cited in~\cite{DS1}.
\bibitem{BudHR} M.~Budden, J.~Hiller and A.~Rapp, Generalized Ramsey theorems for $r$-uniform hypergraphs,
 \emph{Australasian Journal of Combinatorics} 63 (2015) 142--152.
\bibitem{BrBH} M.~Bruce, M.~Budden and J.~Hiller, Lexicographic products of $r$-uniform hypergraphs and some
 applications to hypergraph Ramsey theory, \emph{Australasian Journal of Combinatorics} 70 (2018) 390--401.
\bibitem{Liu} S.~C. Liu, Lower bounds for small Ramsey numbers on hypergraphs, in \emph{Proceedings of COCOON 2019},
 LNCS 11653, Springer (2019) 412--424; arXiv:1906.00132.
\bibitem{MR1} B.~D. McKay and S.~P. Radziszowski, The first classical Ramsey number for hypergraphs is computed,
 \emph{Proceedings of SODA'91}, San Francisco (1991) 304--308.
\bibitem{SV} SeverinVisionary, \emph{hypergraph-ramsey-bounds}, GitHub repository (2026), including issue \#3
 (September 14--16, 2026). \url{https://github.com/SeverinVisionary/hypergraph-ramsey-bounds}.
\bibitem{kissat} A.~Biere et al., \emph{kissat} SAT solver, \url{https://github.com/arminbiere/kissat}.
\end{thebibliography}
\end{document}
